\documentclass[11pt]{article}
\usepackage[margin=1in]{geometry}
\usepackage{amsmath}
\usepackage{booktabs}
\usepackage{enumitem}
\usepackage{hyperref}

\title{\textbf{ADR-01: Business Use Case Selection for Customer Query Analytics}}
\author{Architecture Decision Record}
\date{\today}

\begin{document}

\maketitle

\section{Status}
\textbf{Accepted}

\section{Context}

Customer support operations face critical challenges in ticket routing, annotator quality assessment, and understanding customer pain points. Multiple potential use cases were considered for implementing a comprehensive customer query analytics system:

\subsection{Business Problem}
\begin{itemize}
\item Support tickets frequently misrouted due to unclear classification criteria
\item Annotator performance varies significantly, impacting downstream operations
\item Customer sentiment and emerging issues remain invisible to management
\item Manual routing creates bottlenecks and delays response times
\item No systematic approach to identifying customer risk patterns
\end{itemize}

\subsection{Alternative Use Cases Considered}
\begin{enumerate}
\item \textbf{E-commerce Product Reviews Analysis}: Focus on product feedback and recommendations
\item \textbf{Social Media Sentiment Monitoring}: Track brand perception across social platforms  
\item \textbf{Employee Feedback Analysis}: Internal HR sentiment and engagement tracking
\item \textbf{Customer Support Ticket Classification}: Hierarchical ticket routing and quality assessment
\item \textbf{Sales Call Transcription Analysis}: Lead qualification and sales process optimization
\end{enumerate}

\section{Decision}

We selected \textbf{Customer Support Ticket Classification} as the primary use case for the following reasons:

\subsection{Strategic Alignment}
\begin{itemize}
\item Direct impact on customer experience and satisfaction metrics
\item Clear ROI through reduced response times and improved routing accuracy
\item Addresses immediate operational pain points affecting daily business operations
\item Scalable solution applicable across multiple product lines and support channels
\end{itemize}

\subsection{Technical Feasibility}
\begin{itemize}
\item Well-defined hierarchical classification structure (L1: Technical/Billing/Account, L2: Specific subtypes)
\item Measurable success criteria through inter-annotator agreement analysis
\item Rich dataset enabling comprehensive analytics across multiple dimensions
\item Clear stakeholder requirements from support management, QA teams, and operations
\end{itemize}

\subsection{Business Impact}
\begin{itemize}
\item \textbf{Operational Efficiency}: Automated routing reduces manual triage time by estimated 60\%
\item \textbf{Quality Assurance}: Systematic annotator performance assessment enables targeted training
\item \textbf{Customer Insights}: Sentiment analysis and risk detection provide proactive issue identification
\item \textbf{Strategic Intelligence}: Topic modeling reveals emerging customer needs and product gaps
\end{itemize}

\section{Consequences}

\subsection{Positive Outcomes}
\begin{itemize}
\item \textbf{Immediate Value}: Support team can implement improved routing within weeks
\item \textbf{Measurable Results}: Clear KPIs (routing accuracy, response time, customer satisfaction)
\item \textbf{Scalable Framework}: Architecture extends to other text classification challenges
\item \textbf{Data-Driven Decisions}: Management gains visibility into customer patterns and annotator performance
\end{itemize}

\subsection{Trade-offs and Limitations}
\begin{itemize}
\item \textbf{Domain Specificity}: Solution optimized for support tickets, requiring adaptation for other use cases
\item \textbf{Annotation Dependency}: Quality heavily dependent on consistent labeling policies and annotator training
\item \textbf{Complexity Management}: Multi-dimensional analysis (sentiment, topics, entities) requires careful prioritization
\item \textbf{Change Management}: Success requires buy-in from support staff and process modifications
\end{itemize}

\subsection{Rejected Alternatives and Rationale}
\begin{itemize}
\item \textbf{E-commerce Reviews}: Lower business priority, existing solutions available
\item \textbf{Social Media Monitoring}: Requires different technical stack, less direct ROI
\item \textbf{Employee Feedback}: Sensitive data concerns, different stakeholder requirements
\item \textbf{Sales Call Analysis}: Technology readiness gaps in speech-to-text pipeline
\end{itemize}

\section{Implementation Guidelines}

\subsection{Success Metrics}
\begin{itemize}
\item Krippendorff's alpha $\geq 0.67$ for acceptable annotation reliability
\item Routing accuracy improvement of $\geq 25\%$ over baseline manual classification
\item Customer satisfaction score improvement of $\geq 10\%$ through faster resolution
\item Annotator consistency improvement measured through reduced pairwise disagreement
\end{itemize}

\subsection{Key Stakeholders}
\begin{itemize}
\item \textbf{Primary}: Customer Support Management, QA Teams
\item \textbf{Secondary}: Product Teams, Operations Management  
\item \textbf{Governance}: Data Science Team, IT Architecture
\end{itemize}

\subsection{Critical Dependencies}
\begin{itemize}
\item Consistent labeling policy development and maintenance
\item Annotator training program aligned with classification schema
\item Integration capabilities with existing ticketing systems
\item Change management support for workflow modifications
\end{itemize}

\section{Review and Evolution}

This decision should be reviewed quarterly based on:
\begin{itemize}
\item Business value delivery and KPI achievement
\item Technical performance and scalability requirements
\item Changing business priorities and customer needs
\item Lessons learned from implementation and operation
\end{itemize}

Future evolution may extend the framework to adjacent use cases such as chat transcripts, email inquiries, or product feedback analysis.

\end{document}